\section{A strong converse for repeatable observation processes}
\label{sec:tail62}
The one-continuation testing bound distinguishes the full number of
response types only when the error is below one half. Repeating a
common diagnostic cycle gives a stronger obstruction: a finite hidden
register cannot support more positive-probability observation-tail types
than it has labels. The underlying tail-event principle is classical
\cite{Blackwell62,Cohn62,GS62}. We give the short argument needed here,
including the hidden-emission and nonhomogeneous conventions, before
applying compactness to separately redesigned finite-horizon machines.

\begin{lemma}[Tail capacity at narrow cuts]\label{lem:tailcapacity62}
Let $(S_t,Y_t)_{t\ge0}$ be a hidden emission process with finite
registers and finite output alphabets, whose joint next-state and
emission law depends on the entire past only through $S_t$ and the
external time. Its transition kernels may be nonhomogeneous and
irreversible. Suppose $|S_t|\le k$ at infinitely many deterministic
times. Then there cannot be more than $k$ pairwise disjoint observation
tail events having positive probability. Here an observation tail event
belongs to $\bigcap_{t\ge0}\sigma(Y_{t+1},Y_{t+2},\ldots)$.
\end{lemma}
\begin{proof}
Let $B_1,\ldots,B_r$ be such events. Include the initialization variable,
all past hidden labels, and all past emissions in $\mathcal F_t$.
The Markov property for future emissions gives a version
\[
 \mathbb P(B_i\mid\mathcal F_t)=b_{i,t}(S_t),\qquad
 b_{i,t}(s)=\mathbb P(B_i\mid S_t=s).
\]
Only labels of positive probability are relevant. Since each $B_i$ is
measurable with respect to $\mathcal F_\infty=\sigma(\bigcup_t\mathcal F_t)$,
upward martingale convergence gives
\[
 b_{i,t}(S_t)\longrightarrow\boldsymbol1_{B_i}
                                  \quad\hbox{almost surely and in }L^1.
\]
Consequently, for every sufficiently large $t$ and each $i$, some
positive-probability label $s_i$ satisfies $b_{i,t}(s_i)>1/2$.
There is one common sufficiently large time for the finite collection
of events. Disjointness implies $\sum_i b_{i,t}(s)\le1$, so these
labels are distinct. At an arbitrarily late narrow cut, $r\le k$.
This is the conditional-probability zero--one mechanism of
\cite[Theorem~1 and Section~5.1]{GS62}, applied to the emitted tail;
no stationarity, mixing rate, or lower bound on state masses is used.
\end{proof}

Fix pairwise distinct laws $P_1,\ldots,P_n$ on one finite alphabet
$\Omega$. An $R$-step $k$-label instrument has one initialization row
$\alpha_i$ for each $i$ and common, possibly time-dependent,
nonnegative emission matrices $K_{t,y}$ of size $k\times k$ with
$\sum_yK_{t,y}$ row-stochastic. Let $Q_{i,R}$ be its output law, and set
\begin{equation}\label{eq:producterror62}
 e_{k,R}=\min_{\alpha_i,K_{t,y}}
             \max_{1\le i\le n}\TV(Q_{i,R},P_i^{\otimes R}).
\end{equation}
The minimum is attained in a finite product of compact simplexes.
Initialization may be randomized and may depend on $i$; subsequent
kernels may not depend on an unrecorded value of $i$.

\begin{proposition}[Uniform finite witnesses for product laws]
\label{prop:productconverse62}
If $k<n$, then $e_{k,R}$ is nondecreasing in $R$ and
\begin{equation}\label{eq:productconverse62}
                         \lim_{R\to\infty} e_{k,R}=1.
\end{equation}
In particular, for every $\varepsilon<1$ there is an integer
$R(k,\varepsilon,P_1,\ldots,P_n)$ for which no such instrument has
error at most $\varepsilon$. When the laws and the tolerance are
rational, one can find such an integer by a terminating sequence of
finite real-algebraic feasibility decisions. No effective bound on the
number of decisions, or polynomial-time claim, is made.
\end{proposition}
\begin{proof}
Truncation and contraction of total variation under marginalization
show that $e_{k,R}\le e_{k,R+1}$. Suppose its limit were some $e<1$.
Consider the compact countable product consisting of the initialization
rows and all kernels $(K_{t,y})_{t\ge1}$. For every $R$, impose all
closed constraints
$\TV(Q_{i,r},P_i^{\otimes r})\le e$ for $i\le n,r\le R$.
Each finite system is feasible by attainment in \eqref{eq:producterror62}.
Nested compactness gives one infinite hidden process satisfying every
finite constraint, with common kernels and $k$ labels at every time.

Let $Q_i$ denote its infinite output law and $P_i^\infty$ the product
law. The finite constraints imply $\TV(Q_i,P_i^\infty)\le e$.
Indeed, for any event, approximate its indicator in $L^1$ under the
finite measure $Q_i+P_i^\infty$ by indicators of cylinder events;
these generate the product sigma-field. The bound on event-probability
differences passes to the limit. Equivalently, one can apply martingale
convergence to conditional probabilities under this finite measure.

For each $i$, let $B_i$ be the event that empirical frequencies of all
symbols converge to $P_i$. These are disjoint observation tail events,
and the strong law gives $P_i^\infty(B_i)=1$. Hence
$Q_i(B_i)\ge1-e>0$. Mix the initialization index uniformly, retaining
the same common transition kernels. The resulting hidden process has
output law $Q=n^{-1}\sum_iQ_i$, so $Q(B_i)>0$ for every $i$.
Lemma~\ref{lem:tailcapacity62} contradicts $n>k$. This proves the limit.

For the last assertion, introduce the finitely many stochastic rows as
real variables. Every word probability is a polynomial in their entries;
the targets are rational. Absolute-value slack variables express each
finite total-variation inequality by polynomial equalities and
inequalities. Real quantifier elimination decides feasibility. Test
$R=1,2,\ldots$ until it is infeasible. The limit just proved guarantees
termination. This is a finite-witness search, not an implementation of a
general minimum-width algorithm.
\end{proof}

\begin{lemma}[A diagnostic cycle]\label{lem:diagnosticcycle62}
Suppose the response-type set $\mathcal X$ in
\eqref{eq:futuretype59} is finite, of cardinality $n\ge2$.
There is a finite deterministic control word $B$ of positive length
$\ell$ such that its action on $\mathcal X$ is the identity and its
output laws $P_1,\ldots,P_n$ from the $n$ types are pairwise distinct.
Successive executions of $B$ give independent, identically distributed
$P_i$-blocks from initial type $i$. Any intervening fixed probe controls
may be inserted if their outputs are discarded.
\end{lemma}
\begin{proof}
Every type has an executable seed--prefix representative, as proved in
Theorem~\ref{thm:causalclassification59}. Apply
Lemma~\ref{lem:itinerary59} to all $n$ types with, for example,
$\delta=1/4$. The resulting finite deterministic itinerary $T$ has
pairwise distinct output laws: its disjoint diagnostic events have
probability at least $3/4$ under the corresponding laws. Its command
product induces a permutation $\pi$ of the finite set $\mathcal X$.
Let $d\ge1$ be the order of this permutation and take $B=T^d$.
Its action on types is the identity, and its output laws remain
distinct because their first-$T$ marginals are distinct.

The output law of any fixed continuation depends only on the response
type, not on its chosen representative. Probes emit fresh conditionally
independent observations and have no physical effect. Since $B$ returns
the type deterministically, its conditional output law after any number
of previous blocks is again $P_i$. Induction gives the product law.
Intervening probes also preserve the type; their ignored outcomes do not
change the deterministic controls or the conditional block laws. The
cycle is a return on the \emph{finite response quotient}, not an assumed
identity-product word in the physical group.
\end{proof}

\begin{theorem}[The full-error causal minimum]\label{thm:causalstrong62}
In the repeatable-probe model of Definition~\ref{def:causal59}, let
$n=|\mathcal X|\in\N\cup\{\infty\}$. For every fixed
$0\le\varepsilon<1$,
\begin{equation}\label{eq:causalstrong62}
                 \lim_{N\to\infty}\mathsf{PW}_{N,\varepsilon}=n.
\end{equation}
When $n<\infty$, equality with $n$ holds at every sufficiently large
horizon. For every integer $k<n$ there is $L_{k,\varepsilon}<\infty$
such that every error-$\varepsilon$ realization has
\begin{equation}\label{eq:causalstrongoccupation62}
          \#\{1\le t\le N:|S_t|\le k\}\le L_{k,\varepsilon},
\end{equation}
with no restriction on intervening widths. Thus the finite-response
cardinality is the eventual minimum throughout the whole interval
below the maximal total-variation error, not only below one half.
The full-transcript and fresh nondisturbing probe assumptions are
unchanged. No command return condition is imposed.
\end{theorem}
\begin{proof}
For $n=\infty$, Theorem~\ref{thm:causalclassification59} already proves
\eqref{eq:causalstrongoccupation62} for every $k$ and
$\varepsilon<1$. For $n=1$, nonempty registers and the exact one-type
machine prove the assertion. Suppose $2\le n<\infty$ and $k<n$.
Choose executable representatives of all types, and pad their prefixes
with a fixed probe to one common length $b\ge0$. Let $B$, $\ell$ and
$P_i$ be as in Lemma~\ref{lem:diagnosticcycle62}. Fix an integer $R$
for which $e_{k,R}>\varepsilon$, supplied by
Proposition~\ref{prop:productconverse62}.

Consider an advertised profile. Discard narrow cuts with $t<b$, and
greedily select from the others
$t_1<\cdots<t_J$ with gaps at least $\ell$. If there are $R+1$
selected cuts, start from each of the equal-length representatives,
pad to $t_1$ with the same probe, and put one copy of $B$ immediately
after each of $t_1,\ldots,t_R$. Fill the remaining positions before
the next selected cut with that probe. Complete the horizon with any
fixed controls. For each representative the retained $R$ block outputs
have target law $P_i^{\otimes R}$.

Marginalize all prefix, filler, and post-block outputs. In the proposed
machine, the actual transition from one selected cut to the next,
jointly with the retained block output, is a nonnegative instrument
between at most $k$ labels. Summing over the discarded emissions gives
its stochastic normalization. All its entries are common to the
representatives: the schedule is deterministic and the retained label
is the only private past information available to a row. Intervening
registers may be arbitrarily wide. The cut-$t_1$ distributions are
allowed to be the arbitrary initialization rows $\alpha_i$ in
\eqref{eq:producterror62}. After padding unused labels, these composite
kernels therefore define an $R$-step $k$-label instrument.

Whole-transcript correctness under these deterministic policies and
contraction of TV under marginalization give its error at most
$\varepsilon$. This contradicts $e_{k,R}>\varepsilon$. Thus $J\le R$.
The discarded cuts number at most $b$, and each selected cut accounts
for at most $\ell$ remaining narrow cuts. Consequently one may take
\[
                         L_{k,\varepsilon}=b+\ell R.
\]
This count includes empty eligible intervals and excludes cut zero.
The prefix, cycle, and witness integer depend only on the interface,
$k$ and $\varepsilon$, not on the horizon or machine. No correctness
conditional on a rare observed prefix was used.

An exact stationary $n$-type machine works at every horizon by
Theorem~\ref{thm:causalclassification59}. Applying the occupation bound
to $k=n-1$ shows that a smaller peak is impossible once
$N>L_{n-1,\varepsilon}$. This proves eventual equality and the theorem.
\end{proof}

\begin{corollary}[Strong converse below the response cardinality]
\label{cor:processerrorone62}
If $k<n$, the least worst-policy whole-transcript error achievable by
peak width at most $k$ tends to one as the horizon tends to infinity.
At error one, one label always suffices. When the finite group action
and all probe tables are explicitly rational, the proof supplies a
terminating search for a valid occupation constant via the finite
product-law feasibility problems.
\end{corollary}
\begin{proof}
For each $\varepsilon<1$, the preceding occupation bound excludes a
$k$-peak realization for every $N>L_{k,\varepsilon}$. TV is at most one,
which proves the limit and the endpoint. For explicit finite data,
response types and their command permutations can be enumerated. The
common finite itinerary and its permutation order can be found by finite
search, its cycle laws are rational, and
Proposition~\ref{prop:productconverse62} finds $R$. The claimed bound
$b+\ell R$ then follows. No efficiency conclusion follows from this
termination proof.
\end{proof}

\begin{proposition}[Why one half is not the process boundary]
\label{prop:binarytail62}
For $P_1=\delta_0$, $P_2=\delta_1$ and one hidden label,
\begin{equation}\label{eq:binarytail62}
                         e_{1,R}=1-2^{-R}.
\end{equation}
Nevertheless, a single width-one cut, followed by two-label storage,
can have worst error one half for these two product laws at arbitrarily
large remaining horizons.
\end{proposition}
\begin{proof}
A one-label instrument has independent time-dependent emissions with
probabilities $p_t$ of zero. Its two target-string probabilities are
$u=\prod_t p_t$ and $v=\prod_t(1-p_t)$. The two TV errors are $1-u$
and $1-v$. Since $uv\le4^{-R}$, their maximum is at least $1-2^{-R}$,
attained by $p_t=1/2$ for every $t$. For the last assertion, after the
one-label cut choose a fair bit, store it in two labels, and emit it
repeatedly. Its law is the equal mixture of the two deterministic
strings and is at distance one half from each. Thus for errors above
one half the theorem needs control of \emph{repeated} narrow cuts;
it does not assert the false pointwise lower bound at every interior cut.
\end{proof}

The tail-capacity lemma is a necessary condition for arbitrary hidden
emission processes, including irreversible ones. The diagnostic-cycle
application, however, still uses the stated repeatable-probe model.
It is neither a general HMM/POMDP realization classification nor a
claim about nondisturbing repeated quantum measurements. The new
conclusion is the all-error finite-horizon strong converse and occupation
bound obtained from the classical tail principle, finite-product
compactness, and one common physical diagnostic cycle.
